
\documentclass{article}
\usepackage{colortbl}
\usepackage{makecell}
\usepackage{multirow}
\usepackage{supertabular}

\begin{document}

\newcounter{utterance}

\centering \large Interaction Transcript for game `referencegame', experiment `line\_grids\_rows', episode 8 with reflexion{-}Qwen3.5{-}9B.
\vspace{24pt}

{ \footnotesize  \setcounter{utterance}{1}
\setlength{\tabcolsep}{0pt}
\begin{supertabular}{c@{$\;$}|p{.15\linewidth}@{}p{.15\linewidth}p{.15\linewidth}p{.15\linewidth}p{.15\linewidth}p{.15\linewidth}}
    \# & $\;$A & \multicolumn{4}{c}{Game Master} & $\;\:$B\\
    \hline

    \theutterance \stepcounter{utterance}  
    & & \multicolumn{4}{p{0.6\linewidth}}{
        \cellcolor[rgb]{0.9,0.9,0.9}{
            \makecell[{{p{\linewidth}}}]{
                \texttt{\tiny{[P1$\langle$GM]}}
                \texttt{You have attempted this same task before and did not succeed. These are the plans you wrote after those attempts.} \\
\texttt{Trial \#1: New plan: My previous description was ambiguous because it failed to specify the exact row range for the columns containing "X"s, allowing the distractor grids (which have "X"s in different column patterns or full columns) to potentially match the description if interpreted loosely. To ensure uniqueness, I must explicitly state that the "X"s occupy the first and fourth columns for *all* rows (rows 1 through 5), while the second and third columns are empty for *all* rows. This precise quantification of row coverage will clearly distinguish the target from Distractor 1 (which has "X"s only in the last two columns) and Distractor 2 (which has "X"s in the first two columns).} \\
\\ 
\texttt{Expression: The grid where the first and fourth columns are filled with "X" in every row from 1 to 5, and the second and third columns are empty in every row.} \\
\\ 
\texttt{Trial \#2: New plan: My previous description incorrectly identified the columns containing "X"s. Upon re{-}examining the target grid, the "X"s are located in the **first, fourth, and fifth** columns, not just the first and fourth. The second and third columns are indeed empty. To ensure the expression is accurate and distinguishes the target from Distractor 1 (which has "X"s only in columns 4 and 5) and Distractor 2 (which has "X"s in columns 1 and 2), I must explicitly state that the "X"s occupy the first, fourth, and fifth columns for all rows.} \\
\\ 
\texttt{Expression: The grid where the first, fourth, and fifth columns are filled with "X" in every row from 1 to 5, and the second and third columns are empty in every row.} \\
\\ 
\\ 
\texttt{You are given three grids, where each of them is 5 by 5 in size.} \\
\texttt{Grids have empty cells marked with "▢" and filled cells marked with "X".} \\
\texttt{Your task is to generate a referring expression that best describes the target grid while distinguishing it from the two other distractor grids.} \\
\texttt{The first grid is the target grid, and the following two grids are the distractors.} \\
\\ 
\texttt{Target grid:} \\
\\ 
\texttt{X □ □ X X} \\
\texttt{X □ □ X X} \\
\texttt{X □ □ X X} \\
\texttt{X □ □ X X} \\
\texttt{X □ □ X X} \\
\\ 
\texttt{Distractor grid 1:} \\
\\ 
\texttt{□ □ □ X X} \\
\texttt{□ □ □ X X} \\
\texttt{□ □ □ X X} \\
\texttt{□ □ □ X X} \\
\texttt{□ □ □ X X} \\
\\ 
\texttt{Distractor grid 2:} \\
\\ 
\texttt{X X □ X X} \\
\texttt{X X □ X X} \\
\texttt{X X □ X X} \\
\texttt{X X □ X X} \\
\texttt{X X □ X X} \\
\\ 
\texttt{Instruction: Describe the target grid.} \\
\texttt{Generate the referring expression starting with the tag "Expression: " for the given target grid. Omit any other text.} \\
            }
        }
    }
    & \\ \\

    \theutterance \stepcounter{utterance}  
    & \multicolumn{4}{p{0.6\linewidth}}{
        \cellcolor[rgb]{0.9,0.9,0.9}{
            \makecell[{{p{\linewidth}}}]{
                \texttt{\tiny{[P1$\rangle$GM]}}
                \texttt{Expression: The grid where the first, fourth, and fifth columns are filled with "X" in every row from 1 to 5, and the second and third columns are empty in every row.} \\
            }
        }
    }
    & & \\ \\

    \theutterance \stepcounter{utterance}  
    & & & \multicolumn{2}{p{0.3\linewidth}}{
        \cellcolor[rgb]{0.9,0.9,0.9}{
            \makecell[{{p{\linewidth}}}]{
                \texttt{\tiny{[GM$|$GM]}}
                \texttt{Expression: The grid where the first, fourth, and fifth columns are filled with "X" in every row from 1 to 5, and the second and third columns are empty in every row.} \\
            }
        }
    }
    & & \\ \\

    \theutterance \stepcounter{utterance}  
    & & \multicolumn{4}{p{0.6\linewidth}}{
        \cellcolor[rgb]{0.9,0.9,0.9}{
            \makecell[{{p{\linewidth}}}]{
                \texttt{\tiny{[GM$\rangle$P2]}}
                \texttt{You have attempted this same task before and did not succeed. These are the plans you wrote after those attempts.} \\
\texttt{Trial \#1: New plan: I incorrectly identified the "First" grid because I misread the column positions of the filled cells. The expression requires Xs in the 1st and 4th columns, but the First grid has Xs in the 1st and 2nd columns. The Second grid has Xs in the 4th and 5th columns. The Third grid correctly has Xs in the 1st and 4th columns with empty cells in the 2nd and 3rd. Next time, I will carefully count the column indices for each grid to ensure the filled cells match the specified columns (1 and 4) before selecting. The correct answer is Third.} \\
\\ 
\texttt{Trial \#2: New plan: I failed because I answered "None" instead of selecting one of the required options (first, second, or third). The expression describes a grid with "X" in columns 1, 4, and 5, and empty cells in columns 2 and 3. Comparing this to the grids: the First grid has Xs in columns 1 and 2; the Second grid has Xs in columns 4 and 5; the Third grid has Xs in columns 1, 4, and 5 with empty cells in columns 2 and 3. Therefore, the Third grid matches the description. Next time, I will verify the column patterns against the expression and output the correct ordinal number ("third") as the only text.} \\
\\ 
\\ 
\texttt{You are given three grids, where each of them is 5 by 5 in size.} \\
\texttt{Grids have empty cells marked with "▢" and filled cells marked with "X".} \\
\texttt{You are also given a referring expression that describes one of the given grids.} \\
\texttt{Your task is to select the grid that matches the given referring expression.} \\
\texttt{Generate only the number (in text) of the grid that matches the given expression by selecting first, second, or third.} \\
\\ 
\texttt{First:} \\
\\ 
\texttt{X X □ X X} \\
\texttt{X X □ X X} \\
\texttt{X X □ X X} \\
\texttt{X X □ X X} \\
\texttt{X X □ X X} \\
\\ 
\texttt{Second:} \\
\\ 
\texttt{□ □ □ X X} \\
\texttt{□ □ □ X X} \\
\texttt{□ □ □ X X} \\
\texttt{□ □ □ X X} \\
\texttt{□ □ □ X X} \\
\\ 
\texttt{Third:} \\
\\ 
\texttt{X □ □ X X} \\
\texttt{X □ □ X X} \\
\texttt{X □ □ X X} \\
\texttt{X □ □ X X} \\
\texttt{X □ □ X X} \\
\\ 
\texttt{Expression: The grid where the first, fourth, and fifth columns are filled with "X" in every row from 1 to 5, and the second and third columns are empty in every row.} \\
\texttt{Question: Which grid does the expression refer to?} \\
\texttt{Start with the tag "Answer: ", followed by your selection. Omit any other text.} \\
            }
        }
    }
    & \\ \\

    \theutterance \stepcounter{utterance}  
    & & & \multicolumn{4}{p{0.6\linewidth}}{
        \cellcolor[rgb]{0.9,0.9,0.9}{
            \makecell[{{p{\linewidth}}}]{
                \texttt{\tiny{[GM$\langle$P2]}}
                \texttt{Answer: third} \\
            }
        }
    }
     \\ \\

    \theutterance \stepcounter{utterance}  
    & & & \multicolumn{2}{p{0.3\linewidth}}{
        \cellcolor[rgb]{0.9,0.9,0.9}{
            \makecell[{{p{\linewidth}}}]{
                \texttt{\tiny{[GM$|$GM]}}
                \texttt{third} \\
            }
        }
    }
    & & \\ \\

\end{supertabular}
}

\end{document}
